% Introduction: scoped problem statement and bounded contribution for ORION Paper III.
Scientific claims are distributed across disciplines, each with its own
terminology, measurement practices, contexts, and representational choices.
A global synthesis can therefore fail even when every local statement is
individually well formed: identical surface terms can refer to different
objects, different labels can denote compatible objects, and apparently
contradictory claims can differ only in modality, attribution, or context.
The central problem of this paper is not retrieval itself, but the authority
to identify two already-structured scientific representations as one global
statement.

\subsection{The proposal in brief}

ORION Paper III represents each source as a \emph{source-local scientific
projection} bound to its provenance. A proposed cross-source mapping records
scientific coordinates such as referent, construct, measurement,
context, polarity, modality, and attribution separately rather than reducing
them to one similarity score. When the coordinates required by the claim are
compatible, the projections may be \textsc{glue}d. When a load-bearing
coordinate is incompatible or unresolved, the mapping remains a typed
\emph{obstruction} or plural view rather than being forced into a canonical
merge.

This paper evaluates that mapping rule at the layer for which external
evidence currently exists. The initial public-reference study contains 32
already-structured cases derived from pinned external sources. A second,
disjoint 32-case holdout was selected and execution-frozen before its system
outputs, including its exact gold hash, source revisions, evaluator code,
bootstrap seed, margins, and pass rule. On that confirmatory holdout, ORION
produced zero false merges versus 0.1875 for flat predicate canonicalization;
the paired difference was $-0.1875$ with a 95\% bootstrap interval of
$[-0.34375,-0.0625]$. The false-split difference versus an exact-coordinate
conservative control was $0.000$ with interval $[0.000,0.000]$, satisfying the
predeclared confirmatory rule.

The evidence is intentionally narrower than the full Global Knowledge
Portrait programme. It does not establish raw-text extraction superiority,
expert eight-family construct validity, reader-level recoverability of a
generated portrait, downstream scientific-answer improvement, or necessity
of every semantic coordinate. Those are follow-up claims and remain
\conststatus{} rather than being inferred from the mapping result.

\subsection{Claims and non-claims}

The paper absorbs the following mechanism families as prior work rather than
claiming them independently:

\begin{itemize}
    \item source-grounded scientific structure extraction and knowledge
          assembly (MUSE~\cite{muse2026}, Discovery Engine~\cite{discoveryEngine2025});
    \item multidisciplinary scientific schemas (SciSchema.org~\cite{scischema2026});
    \item evidence-linked schema induction and fusion
          (SCOPE/SCION~\cite{scopes2026});
    \item provenance-aware schema-conditioned integration and retrieval
          (Executable Schema Contracts~\cite{schemaContracts2026});
    \item scientific entity/relation extraction (SciER~\cite{scier2024});
    \item ontology and schema alignment
          (Ontology Matching~\cite{euzenat2007ontology}, LLMATCH~\cite{llmatch2025});
    \item heterogeneous knowledge-graph construction, entity resolution,
          fusion, and provenance~\cite{hofer2024kgconstruction,verma2023scholarlykg};
    \item stance and claim-comparison methods~\cite{hardalov2022stance};
    \item literature-based discovery~\cite{swanson1986,swanson1997interactive};
    \item scientific pluralism~\cite{kellert2006,cartwright1999};
    \item retrieval-augmented scientific synthesis and cross-domain agent
          orchestration~\cite{openScholar2024,biosage2025}.
\end{itemize}

The surviving scoped empirical claim is P3.C5: for the frozen
public-reference mapping task, explicit scientific-coordinate conditions plus
an obstruction state reduce false merges relative to flat predicate
canonicalization while satisfying the predeclared false-split guard. The
initial 32 cases are not pooled into the confirmatory verdict. Zero-effect
coordinate ablations are retained as limitations, not interpreted as proof
that those coordinates are unnecessary.

\subsection{Structure of the paper}

Section~\ref{sec:related} establishes the absorbed prior-work boundary.
Section~\ref{sec:method} defines source projections, mappings, \textsc{glue},
and obstruction. Section~\ref{sec:dataset} describes the two public-reference
mapping datasets and their provenance. Section~\ref{sec:eval} records the
prospective confirmatory design and controls. Section~\ref{sec:results}
reports the frozen result and ablations. Section~\ref{sec:limitations}
separates the supported mapping claim from the broader unexecuted programme,
and Section~\ref{sec:conclusion} states the bounded conclusion.
\endinput
